\section{Related Work}
\label{sec:related}

\subsection{Local and hybrid search for graph coloring}

Tabu search has long been a central approach to graph coloring. Hertz and de Werra introduced a tabu-search method for fixed-$k$ coloring~\cite{hertz1987tabu}, and Tabucol-style local search has since remained an important intensification component in graph-coloring heuristics~\cite{malaguti2010survey}. Many strong methods combine tabu-based local improvement with broader population or memory mechanisms. Representative examples include the hybrid evolutionary algorithm of Galinier and Hao~\cite{galinier1999hybrid}, the adaptive-memory approach of Galinier et al.~\cite{GalinierHertzZufferey2008}, and the HEAD memetic framework~\cite{moalic2018head}. These methods illustrate a recurring design pattern in graph coloring: a broader search mechanism generates or recombines candidate colorings, while local search performs intensive repair. Our terminal RLACO$\rightarrow$Tabucol stage follows this broad complementarity, but differs in how the upstream guidance is obtained: the population-based construction is biased by a frozen RL policy trained in a different search framework, and the transferability of that learned guidance is itself the object of study.

\subsection{Ant colony optimization and learning-enhanced ACO}

ACO constructs solutions incrementally from pheromone and heuristic information. Costa and Hertz adapted this paradigm to graph coloring~\cite{costa1997ants}, establishing ACO as a constructive alternative to single-solution local search. More recent work has incorporated machine learning into ACO in other combinatorial settings. Sun et al.~\cite{sun2022boosting} use learned solution predictions to bias ACO sampling and pheromone initialization, while DeepACO~\cite{ye2023deepaco} uses deep reinforcement learning to learn heuristic guidance for ant-based construction. These approaches demonstrate that learned information can improve ACO, but the learning component is developed for use within the target ACO framework. Our setting asks a different question: whether a policy learned inside RLTabucol can be frozen and reused by ACO without retraining. The proposed soft-injection interface therefore transfers an existing color-selection distribution across search frameworks rather than learning a new ACO-specific policy.

\subsection{Learning-guided combinatorial search and policy transfer}

Machine learning has increasingly been used to guide combinatorial optimization by learning construction policies, search heuristics, and improvement decisions~\cite{bengio2021ml4co,darvariu2024graphrl}. For graph coloring, learning-based approaches include reinforcement-learning and tree-search methods~\cite{huang2019coloring,watkins2023dqn,yang2025mctscoloring}, as well as graph-neural-network models for direct coloring or learned vertex ordering~\cite{wang2024negative,langedal2024gnncoloring}. Hybrid learning--heuristic frameworks also place a learned policy inside a classical search process. RLHO~\cite{cai2019rlho} provides a general RL-guided heuristic-optimization framework, while the RLTCol study~\cite{frijters2023rltabucol} combines learned recoloring decisions with TabuCol for graph coloring.

The cited learning-guided methods primarily evaluate learned guidance within the solver architecture in which it is trained or integrated. Our study focuses instead on cross-framework reuse: whether one frozen policy retains useful information when its role changes from guiding a Tabucol-based pipeline to biasing population-based ACO construction. We therefore distinguish three levels of evidence: usefulness inside the source RLTabucol framework, transferability through the soft pheromone prior, and final optimization value when the best RLACO-generated coloring is passed to a separate downstream Tabucol solver.
